\subsection{Spectrum and resolvent}

Lemma 4. --- Let E be a complex Banach space and let u be a closed injective partial operator on E such that \(u^{-1} \in \mathcal{L}(\mathrm{E})\) . Let \(v \in \mathcal{L}(\mathrm{E})\)

such that \(\|v\| < \|u^{-1}\|^{-1}\) . Then the partial operator \(u + v\) is injective, we have \((u + v)^{-1} \in \mathcal{L}(\mathrm{E})\) and

\[
(u + v) ^ {- 1} = u ^ {- 1} \circ \sum_ {k = 0} ^ {+ \infty} (- v u ^ {- 1}) ^ {k},\tag{6}
\]

where the series is absolutely convergent in \(\mathcal{L}(\mathrm{E})\) . Moreover, we have

\[
\| (u + v) ^ {- 1} \| \leqslant \frac {\| u ^ {- 1} \|}{1 - \| v \| \| u ^ {- 1} \|}.
\]

Since \(\|v\|\|u^{-1}\|<1\) , the series with general term \((-vu^{-1})^{k}\) is absolutely convergent in \(\mathcal{L}(\mathrm{E})\) and its sum is the inverse of the endomorphism \(1_{E}+vu^{-1}\) (prop. 2 of I, p. 22). Consequently, the partial operator \((1_{\mathrm{E}}+vu^{-1})\circ u=u+v\) (IV, p. 226, rem. vi) is injective and the reciprocal partial operator \((u+v)^{-1}=u^{-1}\circ(1_{\mathrm{E}}+vu^{-1})^{-1}\) (IV, p. 226, rem. iii) belongs to \(\mathcal{L}(\mathrm{E})\) . Since

\[
\left\| \left(1 _ {\mathrm{E}} + v u ^ {- 1}\right) ^ {- 1} \right\| \leqslant \sum_ {k = 0} ^ {+ \infty} (\left\| v \right\| \left\| u ^ {- 1} \right\|) ^ {k} = \frac {1}{1 - \left\| v \right\| \left\| u ^ {- 1} \right\|},
\]

the lemma is proved.

DEFINITION 8. --- Let u be a closed operator with dense domain on a complex Banach space E. The resolvent set of u is the set of complex numbers λ such that the partial operator \(\lambda 1_{\mathrm{E}} - u\) is injective and its inverse \((\lambda 1_{\mathrm{E}} - u)^{-1}\) belongs to \(\mathcal{L}(\mathrm{E})\) .

The spectrum of u, denoted \(\mathrm{Sp}(u)\) , is the complement of the resolvent set in C.

If \(\lambda \in \mathbf{C} - \mathrm{Sp}(u)\) , we denote by \(\mathrm{R}(u,\lambda) \in \mathcal{L}(\mathrm{E})\) the inverse of \(\lambda 1_{\mathrm{E}} - u\) . The map from \(\mathbf{C} - \mathrm{Sp}(u)\) into \(\mathcal{L}(\mathrm{E})\) that sends \(\lambda\) to \(\mathrm{R}(u,\lambda)\) is called the resolvent of \(u\) .

Remarks. --- Let E be a complex Banach space and let u be a closed operator with dense domain on E.

\begin{enumerate}
\def\labelenumi{\arabic{enumi})}
\item
  If \(u \in \mathcal{L}(\mathrm{E})\) , its spectrum coincides with the spectrum of the element \(u\) of the algebra \(\mathcal{L}(\mathrm{E})\) (I, p. 2, def. 1).
\item
  Let D be the domain of \(u\) . For every \(\lambda \in \mathbf{C} - \mathrm{Sp}(u)\) , we have
\end{enumerate}

\[
1 _ {\mathrm{E}} = (\lambda 1 _ {\mathrm{E}} - u) \circ \mathrm{R} (u, \lambda), \quad 1 _ {\mathrm{D}} = \mathrm{R} (u, \lambda) \circ (\lambda 1 _ {\mathrm{E}} - u).\tag{7}
\]

Moreover, for \(\lambda_{1}\) and \(\lambda_{2}\) in the resolvent set of u, we have

\[
\mathrm{R} (u, \lambda_ {1}) - \mathrm{R} (u, \lambda_ {2}) = (\lambda_ {2} - \lambda_ {1}) \mathrm{R} (u, \lambda_ {2}) \circ \mathrm{R} (u, \lambda_ {1}),\tag{8}
\]

\[
\mathrm{R} (u, \lambda_ {1}) \circ \mathrm{R} (u, \lambda_ {2}) = \mathrm{R} (u, \lambda_ {2}) \circ \mathrm{R} (u, \lambda_ {1}),\tag{9}
\]

Indeed, one has

\[
\mathrm{R} (u, \lambda_ {1}) - \mathrm{R} (u, \lambda_ {2}) = \mathrm{R} (u, \lambda_ {1}) \circ 1 _ {\mathrm{E}} - 1 _ {\mathrm{D}} \circ \mathrm{R} (u, \lambda_ {2}).
\]

Since \(1_{\mathrm{E}} = (\lambda_21_{\mathrm{E}} - u)\circ \mathrm{R}(u,\lambda_2)\) and \(1_{\mathrm{D}} = \mathrm{R}(u,\lambda_1)\circ (\lambda_11_{\mathrm{E}} - u)\) , we obtain

\[
\begin{array}{c} \mathrm{R} (u, \lambda_ {1}) - \mathrm{R} (u, \lambda_ {2}) = \lambda_ {2} \mathrm{R} (u, \lambda_ {1}) \circ \mathrm{R} (u, \lambda_ {2}) \\ - \mathrm{R} (u, \lambda_ {1}) \circ u \circ \mathrm{R} (u, \lambda_ {2}) - \lambda_ {1} \mathrm{R} (u, \lambda_ {1}) \circ \mathrm{R} (u, \lambda_ {2}) \\ + \mathrm{R} (u, \lambda_ {1}) \circ u \circ \mathrm{R} (u, \lambda_ {2}), \end{array}
\]

whence the first formula. Exchanging the roles of \(\lambda_{1}\) and \(\lambda_{2}\) we deduce the second formula.

\begin{enumerate}
\def\labelenumi{\arabic{enumi})}
\setcounter{enumi}{2}
\item
  Let \(\lambda \in \mathbf{C}\) . By the closed graph theorem (EVT, I, p. 19, cor. 5), if the linear map from \(\mathrm{dom}(u)\) to E defined by \(x \mapsto (\lambda 1_{\mathrm{E}} - u)(x)\) is bijective, then its inverse, whose graph is closed, is continuous from E to E. Therefore \(\lambda\) belongs to the resolvent set of \(u\) if and only if the map \(x \mapsto (\lambda 1_{\mathrm{E}} - u)(x)\) from \(\mathrm{dom}(u)\) into E is bijective.
\item
  If \(\lambda\) belongs to the spectrum of u, one of the following properties holds:
\end{enumerate}

\begin{enumerate}
\def\labelenumi{(\roman{enumi})}
\item
  The kernel of \(\lambda 1_{\mathrm{E}} - u\) is not reduced to 0; one then says that \(\lambda\) is an eigenvalue of \(u\) , and that the dimension of \(\operatorname{Ker}(\lambda 1_{\mathrm{E}} - u)\) is its multiplicity;
\item
  The partial operator \(\lambda 1_{\mathrm{E}} - u\) is injective and its image is not dense in E; one says that \(\lambda\) belongs to the residual spectrum of \(u\) ;
\item
  The partial operator \(\lambda 1_{\mathrm{E}} - u\) is injective, its image is dense in E, but \(\lambda 1_{\mathrm{E}} - u\) is not surjective; one says that \(\lambda\) belongs to the continuous spectrum of \(u\) .
\end{enumerate}

\begin{enumerate}
\def\labelenumi{\arabic{enumi})}
\setcounter{enumi}{4}
\tightlist
\item
  Let \(\lambda\) be a complex number belonging to the resolvent set of \(u\) . The resolvent \(\mathrm{R}(u,\lambda)\) is an injective linear map from E into E; its image is the domain of \(u\) and \(u = \lambda 1_{\mathrm{E}} - \mathrm{R}(u,\lambda)^{-1}\) (cf. IV, p. 226). Conversely, this property characterizes the resolvent set and the resolvent. Precisely, let \(\lambda \in \mathbf{C}\) ; if there exists a continuous injective linear map w from E into E such that \(u = \lambda 1_{E} - w^{-1}\) , then \(\lambda\) belongs to the resolvent set of u and \(w = \mathrm{R}(u, \lambda)\) .
\end{enumerate}

In particular, if v is a closed densely defined operator on E, and if \(\lambda \in C\) is a complex number not belonging to \(\operatorname{Sp}(u) \cup \operatorname{Sp}(v)\) , then the equality \(\operatorname{R}(u, \lambda) = \operatorname{R}(v, \lambda)\) implies that u = v.

\begin{enumerate}
\def\labelenumi{\arabic{enumi})}
\setcounter{enumi}{5}
\tightlist
\item
  We define the spectrum of a closable partial operator \(u\) as the spectrum of its closure.
\end{enumerate}

There exist closed operators whose spectrum is empty, or whose spectrum is equal to C (Exercise 12 of IV, p. 347).

Let E be a complex Banach space and u a closed densely defined operator on E. If F is a complex Banach space and if \(v: E \to F\) is an isomorphism, then we have \(\mathrm{Sp}(v \circ u \circ v^{-1}) = \mathrm{Sp}(u)\) and \(\mathrm{R}(v \circ u \circ v^{-1}, \lambda) = v \circ \mathrm{R}(u, \lambda) \circ v^{-1}\) for all \(\lambda \notin \mathrm{Sp}(u)\) .

PROPOSITION 14. --- Let E be a complex Banach space. Let u be a closed operator with dense domain on E and let U be its resolvent set. a) For every \(\lambda\in\mathrm{U}\) the open disk in C with center \(\lambda\) and radius \(\|\mathrm{R}(u,\lambda)\|^{-1}\) is contained in U;

\begin{enumerate}
\def\labelenumi{\alph{enumi})}
\setcounter{enumi}{1}
\item
  The set U is open in C;
\item
  Suppose that \(\operatorname{Sp}(u)\) is not empty. Let \(\lambda \in \mathrm{U}\) and let \(\delta\) be the distance in \(\mathbf{C}\) of \(\lambda\) to the spectrum of \(u\) . We have \(\delta > 0\) and \(\| \mathrm{R}(u, \lambda) \| \geqslant 1 / \delta\) ;
\item
  The map \(\lambda \mapsto \mathrm{R}(u,\lambda)\) is a holomorphic map from U to \(\mathcal{L}(\mathrm{E})\) . For every integer \(k\in \mathbf{N}\) and every \(\lambda \in \mathrm{U}\) , we have
\end{enumerate}

\[
\frac {\partial^ {k}}{\partial \lambda^ {k}} \mathrm{R} (u, \lambda) = (- 1) ^ {k} k! \mathrm{R} (u, \lambda) ^ {k + 1}.
\]

Let \(\lambda \in \mathrm{U}\) . For every \(\mu \in \mathbf{C}\) such that \(\| (\mu -\lambda)1_{\mathrm{E}}\| < \| \mathrm{R}(u,\lambda)\|^{-1}\) , Lemma 4 applied to the injective partial operator \(\lambda 1_{\mathrm{E}} - u\) and with \(v = (\mu -\lambda)1_{\mathrm{E}}\) implies that the partial operator \(\mu 1_{\mathrm{E}} - u = \lambda 1_{\mathrm{E}} - u + v\) is injective and has a continuous inverse. This implies \(a\) ). By loc. cit., we also have the formula

\[
\mathrm{R} (u, \mu) = \mathrm{R} (u, \lambda) \circ \sum_ {k \in \mathbf {N}} (\lambda - \mu) ^ {k} \mathrm{R} (u, \lambda) ^ {k},
\]

which implies that the resolvent of u is holomorphic in U.

The assertion \(b\) ) follows immediately from \(a\) ). If \(\mathrm{Sp}(u)\) is nonempty, the distance from \(\lambda\) to \(\mathrm{Sp}(u)\) is strictly positive (TG, IX, p. 13, prop. 2), whence \(c\) ).

The last part of the assertion \(d\) ) is proved by induction on \(k\) , the case \(k = 1\) being a consequence of formula (8), p. 245.

PROPOSITION 15. --- Let u be a closed operator with dense domain on a complex Banach space E, and let λ be a complex number belonging to the resolvent set of u.

\begin{enumerate}
\def\labelenumi{\alph{enumi})}
\item
  The subset \(\mathrm{Sp}(\mathrm{R}(u,\lambda)) \setminus \{0\}\) of \(\mathbf{C}\) is the image of the spectrum of \(u\) by the map \(\mu \mapsto (\lambda -\mu)^{-1}\) from \(\mathbf{C} \setminus \{\lambda \}\) into \(\mathbf{C}^*\) ;
\item
  For every \(\mu \neq \lambda\) in \(\mathbf{C}\) , one has
\end{enumerate}

\[
\operatorname{Ker} \left(\mu 1 _ {\mathrm{E}} - u\right) = \operatorname{Ker} \left(\left(\lambda - \mu\right) ^ {- 1} 1 _ {\mathrm{E}} - \mathrm{R} (u, \lambda)\right).
\]

Let us prove assertion \(a\) ). For every \(\mu \neq \lambda\) , we compute

\[
\mu 1 _ {\mathrm{E}} - u = (\lambda - \mu) ((\lambda - \mu) ^ {- 1} 1 _ {\mathrm{E}} - \mathrm{R} (u, \lambda)) (\lambda 1 _ {\mathrm{E}} - u).
\]

Since \(\lambda \notin \mathrm{Sp}(u)\) and \(\mu \neq \lambda\) , the linear map \((\lambda - \mu)(\lambda 1_{\mathrm{E}} - u)\) is a bijection from \(\mathrm{dom}(u)\) onto E. Consequently, this formula implies that \(\mu 1_{\mathrm{E}} - u\) is a bijection from \(\mathrm{dom}(u)\) onto E if and only if \((\lambda - \mu)^{-1} 1_{\mathrm{E}} - \mathrm{R}(u, \lambda)\) is a bijection from E onto E, which implies the assertion.

Let us prove \(b\) ). If \(\mu \neq \lambda\) and \(x \in \mathrm{Ker}((\lambda - \mu)^{-1}1_{\mathrm{E}} - \mathrm{R}(u, \lambda))\) , we have \(x \in \mathrm{dom}(u)\) and the formula \(1_{\mathrm{E}} = (\lambda 1_{\mathrm{E}} - u) \circ \mathrm{R}(u, \lambda)\) implies that \(x \in \mathrm{Ker}(\mu 1_{\mathrm{E}} - u)\) . Conversely, if \(x \in \mathrm{Ker}(\mu 1_{\mathrm{E}} - u)\) and \(\mu \neq \lambda\) , the formula \(1_{\mathrm{dom}(u)} = \mathrm{R}(u, \lambda) \circ (\lambda 1_{\mathrm{E}} - u)\) implies \(\mathrm{R}(u, \lambda)(x) = (\lambda - \mu)^{-1}x\) .

PROPOSITION 16. --- Let u be a closed operator with dense domain on a complex Hilbert space E. The spectrum of \(u^{*}\) is the image of the spectrum of u under complex conjugation and, for every element \(\lambda\) of the resolvent set of u, we have \(\mathrm{R}(u,\lambda)^{*}=\mathrm{R}(u^{*},\overline{\lambda})\) . In particular, if u is self-adjoint, the endomorphism \(\mathrm{R}(u,\lambda)\) is normal for every \(\lambda\notin\mathrm{Sp}(u)\) .

Let \(\lambda\in\mathbf{C} - \mathrm{Sp}(u)\) be an element of the resolvent set of \(u\) . We have \(u=\lambda1_{\mathrm{E}}-\mathrm{R}(u,\lambda)^{-1}\) , hence

\[
u ^ {*} = \overline {{\lambda}} 1 _ {\mathrm{E}} - (\mathrm{R} (u, \lambda) ^ {- 1}) ^ {*} = \overline {{\lambda}} 1 _ {\mathrm{E}} - (\mathrm{R} (u, \lambda) ^ {*}) ^ {- 1}
\]

(IV, p. 236 and prop. 9 of IV, p. 239). By Remark 5, it follows that \(\overline{\lambda} \in \mathbf{C} - \mathrm{Sp}(u^{*})\) and that \(\mathrm{R}(u, \lambda)^{*} = \mathrm{R}(u^{*}, \overline{\lambda})\) . Consequently, the spectrum of \(u^{*}\) is contained in the image of \(\mathrm{Sp}(u)\) by complex conjugation. The equality is obtained by applying this property to \(u^{*}\) since \(u^{**} = u\) (corollary of Proposition 8 in IV, p. 237). The last assertion then follows from formula (9), p. 245.

COROLLARY. --- Let u be a partial self-adjoint operator on a complex Hilbert space E. If E is not zero, then the spectrum of u is not empty.

Suppose that \(\operatorname{Sp}(u)\) is empty. Then u is injective and \(u^{-1} = -\operatorname{R}(u,0)\) is an injective endomorphism of E such that \(\operatorname{Sp}(u^{-1}) \subset \{0\}\) (prop. 15, a)), hence \(\operatorname{Sp}(u^{-1}) = \{0\}\) (I, p. 26, cor. 1). Since \(u^{-1}\) is normal (prop. 16), this implies that \(u^{-1}\) is zero (I, p. 110, example 1), which is a contradiction.

Lemma 5. --- Let E be a complex Hilbert space and let u be a closed densely defined partial operator on E. Let \(\lambda \in \mathbf{C}\) . Suppose there exists a real number \(c > 0\) such that

\[
\| u (x) - \lambda x \| \geqslant c \| x \| \quad \text{for all } x \in \operatorname{dom} (u),\tag{10}
\]

\[
\| u ^ {*} (x) - \overline {{\lambda}} x \| \geqslant c \| x \| \quad \text{for all } x \in \operatorname{dom} (u ^ {*}),\tag{11}
\]

Then \(\lambda\) belongs to the resolvent set of \(u\) and \(\| \mathrm{R}(u,\lambda)\| \leqslant c^{-1}\) .

The hypothesis implies that \(u - \lambda 1_{\mathrm{E}}\) and \(u^{*} - \overline{\lambda} 1_{\mathrm{E}}\) are injective. Let F be the image of \(u - \lambda 1_{\mathrm{E}}\) . The space F is dense in E, since its orthogonal is equal to \(\operatorname{Ker}(u^{*} - \overline{\lambda} 1_{\mathrm{E}})\) (prop. 7, c) of IV, p. 236), which is zero.

Let us prove that the space F is closed. Let \((x_{n})_{n\in\mathbf{N}}\) be a sequence in dom(u) such that the sequence \((u(x_{n})-\lambda x_{n})_{n\in\mathbf{N}}\) converges to \(y\in F\) Inequality (10) implies that the sequence \((x_{n})_{n\in\mathbf{N}}\) is a Cauchy sequence in E. Let \(x\in E\) be its limit. The sequence \((x_{n},u(x_{n}))\) converges to \((x,y+\lambda x)\) in E×E; since the graph of u is closed, we therefore have \(x\in\text{dom}(u)\) and \(u(x)=y+\lambda x\) , which proves that \(y\in F\) .

We conclude that F = E. Thus, the partial operator \(u - \lambda1_{E}\) is bijective, whence \(\lambda \notin \operatorname{Sp}(u)\) (Remark 3). Inequality (10) then implies that \(\|R(u, \lambda)\| \leqslant c^{-1}\) .

PROPOSITION 17. — Let E be a complex Hilbert space and let u be a self-adjoint partial operator on E.

\begin{enumerate}
\def\labelenumi{\alph{enumi})}
\item
  The spectrum of u is contained in \(\mathbf{R}\) ;
\item
  If u is positive, then the spectrum of u is contained in \(\mathbf{R}_{+}\) ;
\item
  Suppose that E is nonzero. Let \(\lambda \notin \operatorname{Sp}(u)\) and let \(\delta > 0\) be the distance from \(\lambda\) to the spectrum of \(u\) . We have \(\| \mathrm{R}(u, \lambda) \| = \delta^{-1}\) .
\end{enumerate}

Let \((a,b)\in\mathbf{R}\times\mathbf{R}\) and \(\lambda=a+ib\) . Let \(x\in\mathrm{dom}(u)\) . Since u is self-adjoint, we have \(\langle x|u(x)\rangle\in\mathbf{R}\) , whence

\[
\begin{array}{r} \| u (x) - \lambda x \| ^ {2} = \| u (x) \| ^ {2} - 2 a \langle x | u (x) \rangle + (a ^ {2} + b ^ {2}) \| x \| ^ {2} \\ = \| u (x) - \overline {{\lambda}} x \| ^ {2}. \end{array}
\]

Suppose that \(b \neq 0\) . We then obtain

\[
\begin{array}{c} \| u (x) - \lambda x \| ^ {2} = \| u (x) - \overline {{\lambda}} x \| ^ {2} \geqslant (\| u (x) \| - a \| x \|) ^ {2} + b ^ {2} \| x \| ^ {2} \\ \geqslant b ^ {2} \| x \| ^ {2}. \end{array}
\]

By Lemma 5, we therefore have \(\lambda \notin \operatorname{Sp}(u)\) , whence assertion \(a\) ).

Suppose that u is also positive. If b = 0 and a \textless{} 0, we obtain similarly for \(x \in \text{dom}(u)\) the inequality

\[
\left\| u (x) - \lambda x \right\| ^ {2} = \left\| u (x) - \overline {{\lambda}} x \right\| ^ {2} \geqslant \left(\left\| u (x) \right\| - a \| x \|\right) ^ {2} \geqslant a ^ {2} \| x \| ^ {2},
\]

hence \(\lambda\notin\mathrm{Sp}(u)\) (loc. cit.), which proves \(b\) ).

Let us finally prove \(c\) ). By prop. 16, the resolvent \(\mathrm{R}(u,\lambda)\) is a normal endomorphism of E. Its norm is therefore equal to its spectral radius (cor. 1 of I, p. 108), whence

\[
\| \mathrm{R} (u, \lambda) \| = \sup _ {\mu \in \operatorname{Sp} (\mathrm{R} (u, \lambda))} | \mu |
\]

(theorem 1 of I, p. 24). The spectrum of R(u,λ) cannot be reduced to \{0\}, for in that case one would have \(\| \mathrm{R}(u,\lambda)\| = 0\), so \(\operatorname{dom}(u)\), the image of R(u,λ), would be zero, and E as well. Prop. 15 therefore implies that

\[
\| \mathrm{R} (u, \lambda) \| = \sup _ {\mu \in \mathrm{Sp} (u)} \frac {1}{| \lambda - \mu |} = \frac {1}{\delta}.
\]

COROLLARY. --- Let u be a self-adjoint operator on E.

\begin{enumerate}
\def\labelenumi{\alph{enumi})}
\item
  The residual spectrum of u is empty;
\item
  For all \(\lambda \neq \mu\) in \(\mathbf{C}\) , the eigenspaces of \(u\) relative to \(\lambda\) and \(\mu\) are orthogonal.
\end{enumerate}

Let us prove a). Let \(\lambda\) belong to the spectrum of \(u\) ; it is a real number (prop. 17). We have \(\mathrm{Ker}(\lambda 1_{\mathrm{E}} - u) = \mathrm{Im}(\lambda 1_{\mathrm{E}} - u)^{\circ}\) (prop. 7 of IV, p. 236), hence \(\lambda\) is not an eigenvalue of \(u\) if the partial operator \(\lambda 1_{\mathrm{E}} - u\) has dense image. This implies by definition that the residual spectrum of \(u\) is empty.

Let us prove \(b\) ). By the proposition, we may assume that \(\lambda\) and \(\mu\) are real. Let \(x \in \operatorname{dom}(u)\) such that \(u(x) = \lambda x\) and \(y \in \operatorname{dom}(u)\) such that \(u(y) = \mu y\) . We then have

\[
\lambda \langle x \mid y \rangle = \langle u (x) \mid y \rangle = \langle x \mid u (y) \rangle = \mu \langle x \mid y \rangle ,
\]

whence \(\langle x \mid y \rangle = 0.\)

Remark. --- If u is a closed symmetric operator that is not self-adjoint, its spectrum is not contained in R (cf. Cor. 10 of IV, p. 257 below), and it is possible that the eigenspaces of u corresponding to λ and \(\overline{\lambda}\) need not be orthogonal (exercise 11 of IV, p. 347).

